% Source: skipgram_negative_sampling.tex
\begin{notechapter}
\chapter{Noise Contrastive Estimation and Negative Sampling}\label{ch:nce-negative-sampling}

\section{Theoretical Foundations}
\subsection{Skip-gram Objective}
The original Skip-gram loss maximizes the log-likelihood 
of context words given targets:

\[
\mathcal{L}_{\text{orig}} = \sum_{(w_t,w_c)} \log P(w_c|w_t)
\]

where \( P(w_c|w_t) = \frac{\exp(\mathbf{v}_{w_t}^\top 
\mathbf{v}_{w_c})}{\sum_{w \in V} \exp(\mathbf{v}_{w_t}^\top 
\mathbf{v}_w)} \).

\textbf{Cross-Entropy Interpretation}: Equivalent to 
categorical cross-entropy between:
\begin{itemize}
    \item True distribution: One-hot over \( w_c \)
    \item Predicted distribution: Softmax over \( V \)
\end{itemize}

\subsection{Noise Contrastive Estimation (NCE)}
Reframes problem as binary classification:
\begin{itemize}
    \item \textbf{Positive}: Observed pairs \( (w_t,w_c) \sim P_d \)
    \item \textbf{Negative}: Sampled pairs \( (w_t,w_i) \sim P_n \)
\end{itemize}

Theoretical NCE objective with exact expectations (to be \emph{maximized}; its negative is the loss):

\[
\mathcal{L}_{\text{NCE}} = \mathbb{E}_{w_c \sim P_d} 
\log \sigma(\Delta_{w_c}) + k \mathbb{E}_{w_i \sim P_n} 
\log \sigma(-\Delta_{w_i})
\]

where \( \Delta_w = \mathbf{v}_{w_t}^\top \mathbf{v}_w 
- \log(kP_n(w)) \).

\subsubsection{Derivation from Binary Cross-Entropy}
For binary labels \( D \in \{0,1\} \):

\[
\mathcal{L} = -D \log \sigma(\Delta) - (1-D) \log \sigma(-\Delta)
\]

Averaging over 1 positive (\( D=1 \)) and \( k \) negatives 
(\( D=0 \)) gives the NCE loss.

\section{Negative Sampling}
\subsection{Practical Simplification}
Negative sampling approximates NCE by:
\begin{enumerate}
\item ignoring \( \log(kP_n(w)) \), i.e.\ \( \Delta_w \approx
\mathbf{v}_{w_t}^\top \mathbf{v}_w \);
\item using Monte Carlo sampling for expectations.
\end{enumerate}

Resulting loss for observed pair \( (w_t,w_c) \) and 
sampled negatives \( \{w_i\}_{i=1}^k \):

\[
\mathcal{L}_{\text{NS}} = \log \sigma(\mathbf{v}_{w_t}^\top 
\mathbf{v}_{w_c}) + \sum_{i=1}^k \log \sigma(-\mathbf{v}_{w_t}^\top 
\mathbf{v}_{w_i})
\]

\subsection{Algorithm Steps}
For each training pair \( (w_t,w_c) \):
\begin{enumerate}
    \item Sample \( k \) negatives \( \{w_i\} \sim P_n(w) \)
    \item Compute embeddings: \( \mathbf{v}_{w_t}, \mathbf{v}_{w_c}, 
    \mathbf{v}_{w_i} \)
    \item Calculate gradients:
    \[
    \nabla_{w_t} = [1-\sigma(p)]\mathbf{v}_{w_c} 
    - \sum_{i=1}^k \sigma(n_i)\mathbf{v}_{w_i}
    \]
    \[
    \nabla_{w_c} = [1-\sigma(p)]\mathbf{v}_{w_t}, \quad 
    \nabla_{w_i} = -\sigma(n_i)\mathbf{v}_{w_t}
    \]
    \item Update embeddings: 
    \( \mathbf{v} \leftarrow \mathbf{v} - \eta \nabla \)
\end{enumerate}

\section{Noise Distribution}
\subsection{Parameterized Form}
Standard implementation uses:

\[
P_n(w) = \frac{\text{count}(w)^{3/4}}{\sum_{w'} 
\text{count}(w')^{3/4}}
\]

\subsubsection{Why \( 3/4 \) Power?}
Balances frequency and diversity:
\begin{itemize}
    \item \textbf{Raw counts (\( \alpha=1 \))}: "the" vs "apple" 
    ratio = \( 10^4 \)
    \item \( \alpha=3/4 \): Ratio = \( (10^4)^{3/4} = 10^3 \)
    \item Prevents oversampling frequent words while maintaining 
    frequency hierarchy
\end{itemize}

\section{Gradient Analysis}
\subsection{Positive Term}
\[
\frac{\partial \mathcal{L}}{\partial \mathbf{v}_{w_t}} 
= [1 - \sigma(p)] \mathbf{v}_{w_c}
\]

\textbf{Behavior} (gradient \emph{ascent} on the objective):
\begin{itemize}
\item \( \sigma(p) \approx 0 \) (poor prediction): large update
toward \( \mathbf{v}_{w_c} \);
\item \( \sigma(p) \approx 1 \) (good prediction): small update.
\end{itemize}

\subsection{Negative Term}
\[
\frac{\partial \mathcal{L}}{\partial \mathbf{v}_{w_t}} 
= -\sum_{i=1}^k \sigma(n_i) \mathbf{v}_{w_i}
\]

\textbf{Behavior}:
\begin{itemize}
\item \( \sigma(n_i) \approx 1 \) (false positive): strong push
away from \( \mathbf{v}_{w_i} \);
\item \( \sigma(n_i) \approx 0 \) (true negative): minimal update.
\end{itemize}

\section{Theoretical Guarantees}
\subsection{Consistency}
NCE is a consistent estimator: as the number of samples grows, maximizing the NCE objective recovers the true conditional probabilities, \( \exp(\mathbf{v}_{w_t}^\top \mathbf{v}_w) \to P(w \mid w_t) \) (the model becomes self-normalized). Negative sampling does \emph{not} have this property for general noise distributions: dropping the \( -\log(kP_n(w)) \) term means that it coincides with NCE only when \( kP_n(w) = 1 \) for all \( w \) (uniform noise and \( k = |V| \)). Increasing \( k \) only makes the Monte Carlo average converge to the expectation (law of large numbers):

\[
\lim_{k \to \infty} \frac{1}{k} \sum_{i=1}^k \log \sigma(-n_i) 
= \mathbb{E}_{w_i \sim P_n}[\log \sigma(-\mathbf{v}_{w_t}^\top 
\mathbf{v}_{w_i})]
\]

\subsection{Embedding Geometry}
Optimal embeddings satisfy:

\[
\mathbf{v}_a^\top \mathbf{v}_b \approx \log \frac{P(a|b)}{kP_n(a)}
\]

This encodes both co-occurrence statistics and noise
correction. With the unigram noise \( P_n(a) = P(a) \), the right-hand side is the pointwise mutual information shifted by \( \log k \): \( \operatorname{PMI}(a,b) - \log k \) (Levy and Goldberg, 2014). Negative sampling therefore learns useful embeddings (a factorization of a shifted PMI matrix) without learning a normalized language model.

\section{Implementation Tradeoffs}
\begin{itemize}
    \item \textbf{Small \( k \)} (5-20): Practical default, 
    balances speed/accuracy
    \item \textbf{Large \( k \)}: Better approximation but 
    \( O(k) \) compute/memory
    \item \textbf{Noise Distribution}: \( P_n(w) \propto 
    \text{count}(w)^{3/4} \) works best empirically
\end{itemize}
\end{notechapter}
